Strong gravitational lensing systems 
are a powerful tool for astrophysics and cosmology.
They have been used to study how dark matter is distributed in galaxies and \rf{galaxy} clusters 
\citep[e.g.,][]{kochanek1991a, blandford1992a, broadhurst2000a,  koopmans2002a, 
bolton2006a, clowe2006a, koopmans2006a, bradac2008a, vegetti2009a, huang2009a, jullo2010a,
tessore2016a, monna2017a, jauzac2018a, 
shajib2019a, meneghetti2020a}.
Furthermore, carefully modeling the mass profiles of galaxy-scale strong lenses for a large number of lensing systems 
%(which can be combined with dynamical mass measurements) 
over a wide range of redshifts makes it possible to  study the structural evolution 
%of the internal structure 
of massive elliptical galaxies \ed{\citep[e.g., ][]{bolton2012b, sonnenfeld2013b, LiR2018a, filipp2023a}}.
They are also uniquely suited to probe 
%\rf{low-mass dark matter halos 
dark matter substructure beyond the local universe and line-of-sight low-mass halos
%in the lensing galaxies and along the line of sight} and 
\citep[e.g.,][]{vegetti2014a, hezaveh2016a, vegetti2018a, ritondale2019a, diazrivero2020a, cagan-sengul2022a, nierenberg2023a}, 
to test the predictions of the cold dark matter (CDM) model.
%%%% Add this here:  galaxy mass profile modeling (gas, star, DM) and evolution (may not mention the added benefit to lensed QSO modeling)
% For redshifts beyond the range that lensing galaxies can typically probe, 
% together with high resolution imaging (from adaptive optics and space telescopes),
% % \HST, the \textit{James Webb Space Telescope}, 
% % or the \emph{Nancy Grace Roman Space Telescope} in the near future), 
% strong lensing as an cosmic telescope magnifies spectral \citep[e.g.,][]{cornachione2018a} 
% and spatial features \citep[e.g.,][]{marshall2007a, patricio2019a, vanzella2020a} of the lensed distant galaxies,
% providing the only way to study the morphology and internal structures of galaxies at sub-kpc scales at high redshifts that can extend to $z > 2$.
%sometimes at physical scales only possible through lensing.} 
%for comparisons with
%sometimes to a physical scale comparable to 
%studies of local galaxies.}


Recent measurements of the Hubble constant \ho span a range of $\sim$10\% \citep[e.g.,][]{abbott2017a, abbott2018b, riess2019a, wong2019a, freedman2019a, freedman2020a, planck2020a, khetan2020a, philcox2020a, choi2020a, dhawan2023a}, 
and significant tension between predictions for \ho based on early-universe observables 
%% don't remove: Planck -- sound horizon; Philcox -- matter-rad equality scale
and direct late-universe measurements remain \citep[e.g.,][]{verde2019a}.
%%%% with due to unidentified systematics or new physics.
%with the value inferred from CMB by assuming \lcdm at the low end and some of the most precise local measurements at the other \citep{riess2019a}.
%  direct measurement from Cepheid-calicbrated SNe~Ia and the value inferred from \lcdm from CMB (riess, Planck).  
%The time delays of multiply-lensed QSOs provide one of the most competitive direct measurement of \ho that have the potential to address this 
%discrepancy \citep[e.g.,][]{wong2019a}}.
Multiply-lensed supernovae (SNe) are 
%potentially superior 
ideal for measuring time-delays and \ho because of their well-characterized light curves,
%on a time scale of months, 
and in the case of Type~Ia, with the added benefit of standardizable luminosity \citep{refsdal1964a, treu2010a, oguri2010a, suyu2023a}.
In recent years, strongly lensed supernovae, 
both core-collapse \citep{kelly2015a, rodney2016a} and Type Ia \citep{quimby2014a, goobar2017a, rodney2021a, chen2022a, goobar2023a, pierel2023a, frye2023a}, have been discovered.
Retrospective searches have found lensed SNe as well, \ed{including} Type Ia's \citep{sheu2023a, magee2023a}.
Time-delay from multiply imaged supernovae  
% together with 
% existing approaches of direct measurement of \ho
% measurements from distance ladders \citep[e.g.,][]{riess2019a, freedman2019a, freedman2020a}
% and lensed quasars with the assumption of flat $\Lambda$CDM \citep[e.g.,][]{suyu2010a, suyu2013a, treu2016a, bonvin2017a, wong2019a} 
can therefore be an important independent way to constrain \ho and address this tension in its measurements \citep[e.g.,][]{goldstein2017a, goldstein2018a, goldstein2019a, wojtak2019a, pierel2019a, suyu2020a}. In fact, \ed{\citet{kelly2023a} and \citet{pascale2024a} reported the first two measurements of \ho using lensed SNe.}
Moreover, time-delay \ho measurements are a powerful complement to other independent measurements of the dark energy equation of state
 %, including BAO, CMB, and SNe~Ia
\citep[e.g.,][]{linder2011a, treu2016a, pierel2021a}.
% , sharma2022a, sharma2023a}.
\ed{Both for cluster/group and galaxy-scale lenses, to systematically discover lensed SNe, we can monitor a sample of strong lenses with the highest estimated supernova rates for the lensed sources \citep[e.g.,][]{shu2018a,shu2021a,suyu2023a}}.

\ed{Static strong lenses can be used to constrain cosmological parameters as well.
\citet{li2024a} showed a large sample, $\mathcal{O}(10^4)$,
in combination with velocity dispersion measurements, can provide competitive constraints, including for the possible evolution of dark energy \citep[e.g.,][]{chevallier2001a, linder2003a}.
Compound lenses, lensing systems with two or more sources at different redshifts, are extremely valuable. 
Even a relatively small sample constitutes a powerful cosmological probe \citep[e.g.,][]{collett2012a, collett2014a, linder2016a, sharma2022a, sharma2023a}. 
Compared with a typical strong lens, with a single source, these are much rarer, and yet there have been some very promising recent discoveries \citep{sheu2024a, dux2024a, bolamperti2024a}.}

 

Finally, for nearby strong lensing galaxies, 
extra-galactic tests of General Relativity can be performed 
by combining lens modeling with 
%rotation curves %or other 
spatially resolved 
%dynamic and
stellar kinematic observations \citep{collett2018a}. 
%\rf{%it would be nice if I can cite tristan, but not necessary
%Such a tension, if not resolved by systematic effects, indicates new physics beyond flat $\Lambda$CDM \citep[e.g.,][]{poulin2019a, agrawal2019a}.
%assuming the Friedmann–Lema\^{i}tre–Robertson–Walker (FLRW) metric, 
%\rf{Among the many significant application of strong gravitational lensing:
%1. curvature; 2. galaxy mass profile modeling (gas, star, DM) and evolution (may not mention the added benefit to lensed QSO modeling); 3 using it as a telescope.
%Remember, I don't have say everything, or cite everything.}
% Beyond the flat $\Lambda$CDM cosmological model, 
% by combining strong lensing time delay of multiply-imaged time-varying sources 
% and SNe~Ia distance measurements, 
% one can 
% determine \ho in a model-independent way and  
% measure the spatial curvature \citep[e.g.,][]{li2018a, taubenberger2019a, collett2019a}, 
% and test the 
% Friedmann–Lema\^{i}tre–Robertson–Walker (FLRW) metric \citep{rasanen2015a}.
% }

\rf{For many of these analyses, 
the available sample size of confirmed strong lenses is a major limiting factor.}
%\rf{Many of the studies above would benefit from a large strong lens sample.}
%The application of strong lensing to cosmology has been limited by the available sample size of the lenses.
In the last few years, several groups have used convolutional neural networks to search for strong lensing systems in photometric surveys, including, in increasing sky coverage, 
%(and descending depth), 
CFHTLS \citep{jacobs2017a}, KiDS \citep{petrillo2017a, petrillo2019a, li2020a}, 
DES \citep{jacobs2019a, jacobs2019b}, and
Pan-STARRS \citep{canameras2020a}.

The DESI Legacy Surveys\footnote{\url{https://www.legacysurvey.org/}}
\citep{dey2019a},
for which at least $z$ band is observed with \ed{4-m telescopes},
covers $\sim$\coverage~deg$^2$,
almost four times the size of the Dark Energy Survey \citep{des2005a} footprint.
We identified over $\sim 3500$ new strong lenses in the Legacy Surveys \citep[][respectively; henceforth, H20, H21, S24]{huang2020a, huang2021a, storfer2024a} by using residual neural networks. 
We have also found 436 lensed quasars \citep{dawes2023a} using autocorrelation for a lensed quasar candidate sample.
The entire catalog of these lens candidates can be found on our project website.\footnote{\url{https://sites.google.com/usfca.edu/neuralens/}}

With so many lens discoveries,
% , and at least an order of magnitude more to be discovered in LSST, Euclid, Roman, to realize the full potential such a large sample of strong lenses, 
a fast and robust
% , fully Bayesian 
modeling pipeline is needed.
% Lenstronomy is an open-source lens modeling code that is widely used. 
% However it is CPU-based and its optimization and sampling approaches can be improved upon. 
We therefore developed \gigal \citep{gu2022a}\ed{: 
a GPU-accelerated}, fully forward-modeling Bayesian pipeline that can speed up the lens modeling time by two orders of magnitude or more. 
Recently, we applied \gigal to \ed{an observed} lensing system, 
DESI-253.2534+26.8843 \citep{cikota2023a}, 
making it the first real lens to be modeled by GPUs.


Here, we present the \projname project. 
With thousands of lens candidates, we can choose the best ones for the purposes of 1) detecting low-mass DM halos and 2) measuring \ho\ --- using the best lensed quasar systems and through a future targeted search for live lensed supernovae among lensing systems that are the most likely to host them.
% So far there are three programs under this project:
% the \hst SNAP Program (Paper~I, or this paper), 
% the DESI \progname (Paper~II, Huang, Inchausti, et al., in prep), 
% and the Keck NIRES Program (Paper~III, Agarwal et al., in prep).
% These systems now need followup observations for confirmation.
For high resolution imaging, 
our \HST SNAP program (ID: 15867, PI: Huang) 
\ed{has observed a subset of the best candidates.}
% has confirmed \emph{all} 51 systems that \ed{have been observed.}  
In this paper --- Paper~I of this series --- 
we describe our \hst program.  
We will also present \ed{the first model using GPUs for a strong lensing system with \hst data.}
% Another SNAP program has also completed its observations (ID: 16773, PI: Glazebrook).
Spectroscopic observations are also being carried out.
% \ed{with the Dark Energy Spectroscopic Experiment \citep[DESI;][]{desi2016a, desi2016b}, Keck, and VLT/MUSE.}
% \citet{tran2022a} used VLT Xshooter and Keck NIR obtained spectra and redshifts
% for lens candidates observed by both \hst SNAP programs.
% The Dark Energy Spectroscopic Experiment \citep[DESI;][]{desi2016a, desi2016b} 
% aims to use the baryon acoustic oscillation (BAO) to measure the effect of dark energy on the expansion of the universe.
% It will obtain optical spectra for 40 million galaxies and quasars, 
% constructing a 3D map spanning the nearby universe to  redshifts $\sim 4$.
% will observe all of our lens candidates.
The first \ed{spectroscopic} results from the Dark Energy Spectroscopic Experiment \citep[DESI;][]{desi2016a, desi2016b} strong lens program  will be presented in Paper~II of this series (Huang, Inchausti et al., in prep).
For 
%a small number of systems 
\ed{$\sim 30\%$ of the systems (preliminary estimate)}, the source redshifts \ed{are too high, such that} the typical emission features (e.g., [\ion{O}{2}]) are beyond the \ed{DESI wavelength} range.
For these, our near-IR spectroscopic results from the Keck~2 Telescope will be reported in Paper~III in this series (Agarwal et al., in prep).


This paper is organized as follows.  
We present our \HST SNAP program in   
\S\,\ref{sec:ls-hst}, followed by a brief description of the 
DESI and Keck spectroscopic observations in \S\,\ref{sec:spect}.
In \S\,\ref{sec:lens-model}, we present the model for one of our \hst systems using \gigal. \ed{Our findings are discussed in \S\,\ref{sec:discussion} and the conclusion is provided in \S\,\ref{sec:conclusion}.}



